\documentclass[11pt]{article}
\usepackage[margin=1in]{geometry}
\usepackage{amsmath,amssymb}
\usepackage[hidelinks]{hyperref}

\title{Rosenberg's conjecture in arXiv:1302.2286 was answered by arXiv:1509.08318}
\author{Literature status packet}
\date{2026-06-21}

\begin{document}
\maketitle

\section*{Status}

\textbf{literature\_already\_answered.}  The source paper records Rosenberg's
conjecture as an open problem.  The conjecture was later explicitly confirmed
by Tikuisis--White--Winter.

\section*{Source passage}

The source/open-problem paper is Ben Hayes,
\emph{An \(l^p\)-Version of von Neumann Dimension for Banach Space
Representations of Sofic Groups II}, arXiv:1302.2286 \cite{Hayes}.

On source PDF page 32, after Theorem 3.4, Hayes writes that by Brown--Ozawa
and Choi--Effros, if every amenable group were in the class \(\mathcal C\), then
\(C^*_\lambda(\Gamma)\) would be quasidiagonal for every amenable group.  The
paper then says that this is Rosenberg's conjecture and that it is an open
problem of major current interest.

\section*{Supporting answer}

The supporting paper is Aaron Tikuisis, Stuart White, and Wilhelm Winter,
\emph{Quasidiagonality of nuclear C*-algebras}, arXiv:1509.08318 \cite{TWW}.
Its abstract states as a consequence of the main theorem that it confirms the
Rosenberg conjecture: discrete amenable groups have quasidiagonal
\(C^*\)-algebras.

This directly answers the source status remark: for every discrete amenable
group \(\Gamma\), the reduced group algebra \(C^*_\lambda(\Gamma)\) is
quasidiagonal.

\section*{Scope limitations}

This is a literature-status packet only.  It does not answer Hayes's closing
Conjectures 4.1--4.3 about \(l^p\)-dimension.  It records the later resolution
of the Rosenberg-conjecture open problem explicitly cited in the source paper.

\begin{thebibliography}{9}
\bibitem{Hayes}
B. Hayes,
\emph{An \(l^p\)-Version of von Neumann Dimension for Banach Space
Representations of Sofic Groups II},
arXiv:1302.2286, 2013.

\bibitem{TWW}
A. Tikuisis, S. White, and W. Winter,
\emph{Quasidiagonality of nuclear C*-algebras},
arXiv:1509.08318, 2015.
\end{thebibliography}

\end{document}
